\section*{Step 284: Conrey 1989 Falsification Attempt}

\paragraph{Target.}
The step tests whether Conrey's positive-proportion theorem gives a bridge from more than two fifths of the zeros on the critical line to all zeros on the critical line.

\paragraph{Published theorem.}
Conrey proves that at least \(2/5\) of the zeros of the Riemann zeta-function are simple and on the critical line.  The method is a refinement of Levinson's method and uses a longer mollifier.

\paragraph{Cascade need.}
RH closure requires
\[
  \frac{N_0(T)}{N(T)} \to 1
\]
and, more strongly, exclusion of every off-line nontrivial zero.  A lower bound such as \(N_0(T)/N(T) \ge 2/5\) leaves the remaining zeros uncontrolled.

\paragraph{Obstruction.}
The Conrey--Levinson method gives lower bounds on critical-line proportion through mollified second moments.  Later refinements such as Bui--Conrey--Young move the constant to \(0.4105\), not to \(1\).  The audited method lineage is therefore a positive-proportion apparatus, not a closure apparatus.

\paragraph{Verdict.}
\[
  \boxed{V_{\mathrm{conrey\_RH\_bridge\_blocked}}}
\]
The falsification attempt fails: Conrey 1989 does not provide a 100\% bridge.  The empirical downgrade pattern becomes 8-of-8.
